% ======================================================================
% Section — The scaling-closure lemma: derivation and calibration
% ======================================================================

\section{The scaling-closure lemma: derivation and calibration}\label{sec:scale}

This section is the derivation record behind Theorem~\ref{thm:SC}. Its
order is the order of the reduction: the cascade model and the
dilate-intersection step (\S\ref{sec:scale-model}); the pool law, its
saturation tables, and the status of the $840$ identity
(\S\ref{sec:scale-pool}); the cut law---spacing, $k$-cancellation, and
the conditional comparison that fixes the depth formula's two ends
(\S\ref{sec:scale-cut}); the two-regime structure of the dilate cuts,
with the reflection fiber and the chain class proved
(\S\ref{sec:scale-h2}); the volume law at $T=8$ with its
threshold form, the renormalization lemma, the character system, the
reduction of the remaining obligation to the strand
census, and the transfer-matrix session that closed the third route by
theorem (\S\ref{sec:scale-h1}); and the epistemic status
(\S\ref{sec:scale-status}). Throughout, a claim is labeled
\emph{proved} (with its verification record), \emph{measured} (with its
scope, sample size, and, where sampled, an interval), or \emph{fit}
(with the data it is fit to); extrapolations beyond the tested range
are labeled \emph{projected}; and no claim is made about the number of
sessions or amount of work needed to close the open obligations.

\subsection{The model: the cascade as a dilate-intersection}\label{sec:scale-model}

Fix the rung $T\ge 8$, the prime $p=Tk+\epsilon$, and the composite
modulus $N=p^2$. The frontier cell $(1K,mU)$, $m=T-3$, consists of the
covering families with one kernel member and $m$ unit members; by
Lemma~\ref{lem:rungT} the kernel covers the fibers of a kernel bad set
$B_k$ of size $2k+1$, leaving $|U|=(T-2)k+\epsilon-1$ uncovered fibers
$F_j\cong\Z_p$, and on each fiber the $m$ unit footprints are arithmetic
progressions whose sizes lie in the two-value window and whose starts
obey the drift law. In the census normal form---the normalizer arc
pinned with difference $1$ at the origin, the remaining differences
canonical in $[2,D]$, $D=(p-1)/2$---the placement on fiber $j$ is the
relative-start tuple
\[
c(j)\;=\;c^0(j)-\mu\,j\;\in\;(\Z_p)^{m-1},
\]
where $c^0(j)$ is the residue-determined part (a function of the unit
residues and the three-tier size assignment $\mathrm{sz}(j)$) and
$\mu\in(\Z_p)^{m-1}$ is the free drift, one coordinate per moving arc,
determined by the lifts. The fiber is covered if and only if
$c(j)\in\mathrm{Sol}(\mathrm{sz}(j))$, the census solution set of the
size tuple and the family's difference tuple; this is exactly the
committed ground-truth check. The \emph{fiber cascade} processes the
uncovered fibers one by one, retaining the anchor placements that
survive each; the family dies when the survivor set empties, and the
\emph{cascade depth} is the number of fibers needed.

\begin{lemma}[two-fiber reduction]\label{lem:twofiber}
Fix an anchor fiber $j_0\in U$ and an anchor solution
$a\in\mathrm{Sol}(\mathrm{sz}(j_0))$. The anchor determines the drift,
$\mu=(c^0(j_0)-a)\,j_0^{-1}$. For a second fiber $j$ with
$\lambda:=j\,j_0^{-1}\bmod p$, the anchor solution $a$ survives fiber
$j$ if and only if
\[
\lambda\,a+\tau\;\in\;\mathrm{Sol}(\mathrm{sz}(j)),
\qquad
\tau=c^0(j)-\lambda\,c^0(j_0)\;\in\;(\Z_p)^{m-1}.
\]
Hence the two-fiber candidate set is the dilate-intersection
$\mathrm{Sol}_0\cap\lambda^{-1}(\mathrm{Sol}_j-\tau)$, and the cascade
after $i$ processed fibers is the iterated intersection
$\bigcap_{\ell\le i}\lambda_\ell^{-1}(\mathrm{Sol}_\ell-\tau_\ell)$.
\end{lemma}

\begin{proof}
Substitute $\mu=(c^0(j_0)-a)\,j_0^{-1}$ into the covering condition
$c^0(j)-\mu\,j\in\mathrm{Sol}(\mathrm{sz}(j))$ and rearrange; all maps
are affine over $\Z_p$.
\end{proof}

The cascade step is thus exactly a \emph{dilate-intersection of census
solution sets}: the covering constraints across fibers are affine
conditions on the anchor placement, with dilation ratios
$\lambda=j/j_0$ given by the fiber set. This is the precise form of the
picture under which the anchored-core analysis was organized. Write
$\bar\lambda=\min(\lambda,p-\lambda)$ for the cyclic representative.
The dilate $\nu\mapsto\lambda\nu$ maps the junction range into itself
strongly when $\lambda$ is a \emph{small rational}
($\bar\lambda\le\Lambda_0$, or $\lambda\equiv a/b$ with $a,b\le O(L)$
for the junction length $L$): these are the \emph{resonant} fibers,
where the cut is weak. The remaining fibers impose genuine spacing
(Lemma~\ref{lem:spacing}). At the committed zoo prime ($p=17$, $k=2$)
the resonance zone $\bar\lambda\le 4$ covers half the dilations---the
small-$p$ extreme; its fraction vanishes as $p$ grows. The per-pair
\emph{cut quality} is measured by
\[
\gamma\;=\;\frac{\#\bigl(\mathrm{Sol}_0\cap\lambda^{-1}
(\mathrm{Sol}_j-\tau)\bigr)}
{|\mathrm{Sol}_0|\,|\mathrm{Sol}_j|\,/p^{m-1}},
\]
the survivor count against the random-placement baseline; $\gamma=0$ is
an \emph{instant kill} (the actual intersection is smaller than the
random model), $\gamma\approx 1$ is no cut at all, and $\gamma\gg 1$
marks a resonant pair whose solution sets are too small for the random
baseline to be meaningful.

\subsection{The pool law}\label{sec:scale-pool}

\begin{lemma}[pool law]\label{lem:pool}
The family pool after the $\pm$-collision reduction is bounded by the
distinct canonical difference tuples:
\[
P(T,p)\;\le\;(D-1)(D-2)\cdots\bigl(D-(T-4)\bigr)
\;\le\;\bigl((p-1)/2\bigr)^{T-4},
\qquad D=(p-1)/2,
\]
with difference count $T-4$ (the normalizer's difference is pinned to
$1$).
\end{lemma}

\begin{proof}
The census key of a $\pm$-distinct family is an ordered tuple of
distinct canonical differences in $[2,D]\subset[1,D]$; the falling
factorial counts them.
\end{proof}

The law that is \emph{not} a theorem is the saturation identity
$P=(D-1)_{T-4}$. The committed zoo point ($T=8$, $(1K,5U)$, $p=17$,
$k=2$) records exactly $840$ $\pm$-distinct families and
$840=7\cdot6\cdot5\cdot4$---the bound saturated, every distinct tuple
admitting a covering placement. Table~\ref{tab:pool} is the verdict at
the demanded primes, and it is unambiguous: \emph{the identity is a
local accident of the slackiest prime, and the pool law is
rung-monotone}.

\begin{table}[htbp]
\centering
\small
\caption{Pool saturation at the frontier cells. The $T=8$ rows are
exhaustive censuses (two independent decision procedures, cross-checked
against a third brute force; \S\ref{sec:scale-sound}). The $T=9$ rows
are stride-sampled with full decisions per sampled tuple: saturation
carries a $95\%$ Wilson interval with the stated $n$; the projected
pools (bound $\times$ sampled fraction, one prime each) are
\emph{projected}, not established.}
\label{tab:pool}
\begin{tabular*}{\textwidth}{@{\extracolsep{\fill}} l r l r l l}
\toprule
rung & $p$ & $(k,\epsilon)$ & bound $(D-1)_{T-4}$ & admissible &
saturation \\
\midrule
$T=8$ & $17$ & $(2,1)$ & $840$ & $840$ & $1.000$ \\
$T=8$ & $19$ & $(2,3)$ & $1{,}680$ & $960$ & $0.571$ \\
$T=8$ & $23$ & $(2,7)$ & $5{,}040$ & $1{,}800$ & $0.357$ \\
$T=8$ & $29$ & $(3,5)$ & $17{,}160$ & $2{,}280$ & $0.133$ \\
$T=8$ & $31$ & $(3,7)$ & $24{,}024$ & $360$ & $0.015$ \\
\midrule
$T=9$ & $29$ & $(3,2)$ & $154{,}440$ & --- &
$0.903\ [0.889,0.915]$ \\
$T=9$ & $31$ & $(3,4)$ & $240{,}240$ & --- &
$0.490\ [0.468,0.511]$ \\
\bottomrule
\end{tabular*}
\end{table}

Three readings. \emph{First}, the saturation mechanism: a covering by
arcs with pairwise-distinct differences must waste overlap---the chain
mechanisms of the repeated-difference core all exploit pairing (the
interval tiling needs $d=1$ twice, the parity mechanism needs
$d=d'$, the antipodal needs the $q$-pairs), so saturation requires the
overlap budget to absorb the \emph{distinctness penalty}. The
budget-versus-penalty race is lost at $T=7$ (the distinct pool is empty
for $k\ge 2$: the committed $18$-prime census records $24/24$ at $p=11$
and $36/60$ at $p=13$ at $k=1$, and zero at every $k\ge2$ prime),
marginal at $T=8$ (it holds only at $p=17$, the prime with the largest
overlap window, $\mathrm{ov}\in[3,8]$; it collapses as the window
tightens and as $k$ grows), and won with margin at $T\ge 9$ (at $T=9$
the extra arc's budget $\mathrm{ov}\approx(T-6)k+O(1)$ covers the
penalty for most tuples: $90.3\%$ at the slackier $\epsilon=2$ window,
$49.0\%$ at the tighter $\epsilon=4$). \emph{Second}, the projected
$T=9$ pools: $1.39\times10^5$ (interval $[1.37,1.41]\times10^5$) at
$p=29$ and $1.18\times10^5$ ($[1.12,1.23]\times10^5$) at $p=31$---the
measured $\sim\!10^5$ pool scale at $T=9$ is (approximately) the
distinct-tuple pool, and the $T=10$ reading ($\sim\!10^8$ at $k\ge4$,
the falling factorial times the anchor multiplicity) is consistent with
the rung-monotone law; the $T=10$ saturation itself (predicted
$\ge 0.9$ at the $40\%$ budget) is \emph{not measured}, and at
$k\ge4$ the pools exceed the cascade's reach in any case.
\emph{Third}, the set-form: admissibility is invariant under permuting
the differences ($0$ order-conflicts over the $6{,}720$ checked tuples
at $p=19$ and $p=23$, forced by the size-assignment symmetry of the
census), so the pool is $\#\text{sets}\times(m-1)!$ with
$\#\text{sets}\le\binom{D-1}{m-1}$, and the saturation statement
becomes ``all $\binom{D-1}{m-1}$ sets admissible''---true only at the
slackiest (rung, prime) pairs, asymptotically in $T$. The absolute
admissible count at $T=8$ grows across the table
($840\to960\to1{,}800\to2{,}280$) but collapses at $p=31$ ($360$): the
pool is real but $\epsilon$-modulated, never a smooth polynomial in
$p$.

\subsection{The cut law: spacing, \texorpdfstring{$k$}{k}-cancellation, and the comparison}\label{sec:scale-cut}

\begin{lemma}[dilate-count spacing bound]\label{lem:spacing}
Let $p$ be prime, $\lambda\in\Z_p^*$, $0\le L<p$, $t\in\Z_p$. Then
\[
N(\lambda,t)\;:=\;\#\{\nu\in\{0,\dots,L\}:
(\lambda\nu+t\bmod p)\in\{0,\dots,L\}\;\}
\;\le\;\lfloor L/b_{\min}\rfloor+1,
\]
where $b_{\min}(\lambda,L):=\min\{b\ge 1:\ \|\lambda b\|_p\le L\}$ is
the first return of the $\lambda$-walk to the $L$-ball ($b_{\min}=L+1$
if no such $b\le L$ exists, giving $N\le 1$).
\end{lemma}

\begin{proof}
Let $\nu_1<\nu_2$ be two solutions. Then
$\delta=\nu_2-\nu_1\in[1,L]$ satisfies
$\lambda\delta\equiv(\lambda\nu_2+t)-(\lambda\nu_1+t)$ with both
right-hand residues in $[0,L]$, so $\|\lambda\delta\|_p\le L$; hence
$\delta\ge b_{\min}$. The solutions are therefore $b_{\min}$-separated
in $\nu$, and at most $\lfloor L/b_{\min}\rfloor+1$ of them fit in
$[0,L]$.
\end{proof}

The corollary is the resonant/non-resonant split. The survival fraction
per coordinate is
$\sigma(\lambda):=\max_t N(\lambda,t)/(L+1)\le 1/b_{\min}+1/(L+1)$:
for an integer dilate $\bar\lambda\in[2,L]$ one has $b_{\min}=1$ and
$N\approx(L+1)/\bar\lambda$, so the cut is weak for $\bar\lambda\le 2$
and real for $\bar\lambda\ge 3$; for a small rational
$\lambda\equiv a/b$ with $a,b\le O(L)$ one has $b_{\min}\le b$ and the
cut degrades by the resonance factor, the number of such dilations among
the fibers being $\le\sum_{b\le B}(2L+1)$ for strength $B$---the
resonance budget; and otherwise (no good small-rational representative
---the bulk of the fibers at large $p$) $b_{\min}\gtrsim p/(2L+1)$, so
$\sigma(\lambda)\lesssim(2L+1)/p+1/(L+1)$, and with $L=Rk$, $p=Tk$:
$\sigma\lesssim 2R/T$. One regime caveat must be recorded: the box form
requires $L=Rk<p/2$, i.e.\ $R<T/2$---true for the thin majority of zoo
families ($R\approx4.8$ at $T=8$) but \emph{false} for the fat top
families ($R$ up to $\approx10.5$, $k$-free). For those, the cut is
carried not by the spacing lemma but by the solution sets' own sparsity
(their volume is far below $p^{m-1}$; the measured instant kills),
which is exactly H2m's content.

\begin{proposition}[$k$-cancellation]\label{prop:kcanc}
Assume the volume law in either proved or hypothesized form---for the
chain class, $|\mathrm{Sol}|\approx\mathrm{ov}^{\,m-1}$ with
$\mathrm{ov}\approx(T-6)k$ (Lemma~\ref{lem:chainbox}); for general
families under H1g, $|\mathrm{Sol}|\le c_1\,\mathrm{ov}^{\,m-1}$. Then
the per-fiber cut ratio at the frontier cell is $k$-free:
\[
\mathrm{cut}\;\approx\;\frac{p^{m-1}}{|\mathrm{Sol}|}
\;\gtrsim\;\Bigl(\frac{T}{T-6}\bigr)^{m-1}\big/\bigl(\text{order and
tie factors}\bigr).
\]
\end{proposition}

\begin{proof}
The junction volume and the modulus scale together:
$p^{m-1}=(Tk)^{m-1}$ against $|\mathrm{Sol}|\approx((T-6)k)^{m-1}$,
so the $k^{m-1}$ in the denominator cancels the $k^{m-1}$ in the
numerator and the ratio is a function of $T$ alone, up to the order
factors (the $(m-1)!$ orders and the tie bookkeeping), which are
$k$-free by the census rigidity.
\end{proof}

The earlier form of this proposition was stated with a fitted junction
scale $R$ ($|\mathrm{Sol}|\le(Rk)^{m-2}$, $R\approx4.2$); the chain-class
analysis of \S\ref{sec:scale-h2} falsified that exponent for the
dominant family (the correct chain exponent is $m-1$, not $m-2$) and
\emph{derived} the volume constant from the overlap budget $T-6$. The
fitted $R$ is retired; under H1g the thin families' constant is $c_1$,
bounded by data. Either form leaves the cut $k$-free---the cancellation
is robust to the correction.

\begin{proposition}[comparison; the independent-extension model,
evaluated under the calibrated constants]
\label{prop:depth}
\emph{Model statement, not a cascade theorem}: in the
independent-extension composition---the idealization in which the
per-fiber cuts act on independent copies rather than on the actual
history-dependent survivor sets---the candidate count after $i$
processed fibers is
$C_i\approx\bar\gamma^{\,i}\,|\mathrm{Sol}|^{i+1}/p^{(m-1)i}$. The
actual-cascade statement is Theorem~\ref{thm:SC}'s clause (i) with
Proposition~\ref{prop:Hq}'s induction, conditional on the quantified
ledger; the model below is its calibration. In the model, the family
dies within
\[
D\;\le\;\frac{\log\!\bigl(c_1\,\mathrm{ov}^{\,m-1}\bigr)}
{(m-1)\log\!\bigl(T/(T-6)\bigr)-\log c_1}
\;+\;N_{\mathrm{res}}\;+\;O(1)
\]
fibers, $N_{\mathrm{res}}=1+2\Lambda_0(T-2)/T$. Consequently:
\emph{(a)} (never super-linear) $D\le(T-4)+N_{\mathrm{res}}+O(1)\le
T+3$---the linear law with slope $1$ is the worst case;
\emph{(b)} ($T$-flat at practical $k$) the pool-route term stays in
$[1.1,2.9]$ for $k\le10$, $8\le T\le20$, so the typical depth is
$\approx 2$--$3$ plus the tail, $\approx 9$--$12$;
\emph{(c)} ($k$-logarithmic at fixed $T$) the leading term approaches
$T-4$ only at $k\gtrsim(T/(T-6))^{T-4}$---beyond any prime where the
cascade runs;
\emph{(d)} (closure margin) the worst case $(T-4)+N_{\mathrm{res}}<
|U|=(T-2)k$ for every $T\ge 8$, $k\ge2$ (at $k=2$: $T+3<2T-4\iff
T>7$), with margin growing linearly in $k$; the $\gamma$-tolerance for
closure is $\bar\gamma$ up to $\sim Tk\,(T/(T-6))^{T-4}$, against a
measured effective $\bar\gamma\le\sim\!10$ at the committed cells.
\end{proposition}

\begin{table}[htbp]
\centering
\small
\caption{The depth calibration at the three rungs. Measured: the
committed cascade record (maxima over $k\le9/8/3$; medians $2/3/5$).
Formula: Proposition~\ref{prop:depth} evaluated with derived constants
($c_1=1$, $\Lambda_0=4$) at the top-of-window overlap; the window
shifts the leading term by $\pm1.5$ at $T=8$. Worst case: the
$(T-4)+N_{\mathrm{res}}$ bound. The $T=9$ excess ($11$ measured against
$9.2$) is within the $\gamma$-tolerance and is localized to
$\bar\gamma(T)$ or the near-tiling contamination of the $T=10$ median
($5$ against $2.7$--$4.7$)---decidable only by H2m, not by further
measurement.}
\label{tab:depthcal}
\begin{tabular*}{0.9\textwidth}{@{\extracolsep{\fill}} r r r r r}
\toprule
$T$ & $k$ & measured max & formula & worst case \\
\midrule
$8$ & $2$ & $9$ & $8.5$ & $\le 11$ \\
$9$ & $3$ & $11$ & $9.2$ & $\le 12.2$ \\
$10$ & $3$ & $10$ & $10.1$ & $\le 13.4$ \\
\bottomrule
\end{tabular*}
\end{table}

The calibration carries its provenance explicitly: the per-fiber cut
ratios ($\approx120/960/2{,}160$), the medians, and the maxima of
Table~\ref{tab:depthcal} are the program's three-rung cascade record
from the $T=8/9/10$ measurement runs; those runs predate the present
artifact bundle, whose committed cascade scripts cover the committed
cells only (depths median $2$, maximum $4$; \S\ref{sec:scale-sound}).
The three-rung record is therefore \emph{calibration}, not proof input,
and the derivation of the model rests on the committed machinery
(Lemmas~\ref{lem:twofiber} and \ref{lem:rungT}) alone. Three rungs
cannot distinguish a real growth of $\bar\gamma(T)$ from the
measurement contamination already flagged at $T=10$; what the mechanism
adds is the \emph{localization}---the growth, if real, lives in
$\bar\gamma$ (the equidistribution quality of the junction draws across
fiber dilations) or in the junction fatness, not in the
pool-versus-cut race, whose terms are pinned by
Lemmas~\ref{lem:pool}--\ref{prop:kcanc}.

\subsection{The two-regime structure: the measured \texorpdfstring{$\gamma$}{gamma}, the reflection fiber, and the chain class}\label{sec:scale-h2}

The zone-resolved measurement of the cut quality $\gamma$ over the
committed cells ($764$ admissible fiber pairs; five cells; canonical
signs; the identical seed reproducing the zoo random-$40$) splits into
two regimes, and the split is sharp at every cell.

\begin{table}[htbp]
\centering
\small
\caption{Zone-resolved cut quality $\gamma$ (survivors against the
random baseline) over the committed cells: pair counts, median, and
maximum per zone. Mid zone: $\bar\lambda>4$ and $\bar\lambda<p-4$;
small zone: $\bar\lambda\le 4$. The tail (up to $2.09\times10^4$)
lives entirely in the small zone; its dominant component at
$T=8$-critical is the $\lambda=-1$ reflection fiber of
Lemma~\ref{lem:reflect} (survival fraction exactly $1$).}
\label{tab:gamma}
\begin{tabular*}{\textwidth}{@{\extracolsep{\fill}} l r r r r r r}
\toprule
& \multicolumn{3}{c}{mid zone ($\bar\lambda>4$)}
& \multicolumn{3}{c}{small zone ($\bar\lambda\le4$)} \\
\cmidrule(lr){2-4}\cmidrule(lr){5-7}
cell & $n$ & median & max & $n$ & median & max \\
\midrule
$T=7$, $p=17$, top-$10$ & $40$ & $0.0$ & $0.0$ & $68$ & $0.0$ & $204.7$ \\
$T=7$, $p=29$, top-$10$ & $120$ & $0.0$ & $40.6$ & $70$ & $0.0$ & $508.1$ \\
$T=8$-crit., $p=17$, top-$10$ & $4$ & $0.0$ & $0.0$ & $11$ & $409.4$ & $818.8$ \\
$T=8$ zoo, $p=17$, top-$10$ & $40$ & $0.0$ & $10.7$ & $70$ & $0.0$ & $193.3$ \\
$T=8$ zoo, $p=17$, random-$40$ & $128$ & $0.0$ & $17.6$ & $213$ & $0.0$ & $2.09{\times}10^4$ \\
\bottomrule
\end{tabular*}
\end{table}

The reading is threefold. \emph{(i)} The non-resonant zone satisfies
H2m with enormous margin: the median pair is an instant kill---the
actual intersection is \emph{smaller} than the random model, which is
the lattice sparsity---and the worst mid-zone excess over all five
cells is $\gamma=40.6$ (a single pair at $T=7$, $p=29$); a uniform
$\gamma_0\approx40$ covers every committed mid-zone pair, with
$90$th percentiles in $14$--$422$. \emph{(ii)} The $\gamma$ tail is
entirely a $\bar\lambda\le4$ phenomenon. \emph{(iii)} The tail's
dominant component is proved structure, not noise: at $T=8$-critical
every small-zone $\gamma\in[205,819]$ is a $\lambda=-1$ fiber with
survival fraction exactly $1$ (the random baseline is tiny because the
solution sets there are nearly rigid), and the extreme zoo values
($\le2.09\times10^4$) occur at pairs with
$|\mathrm{Sol}_j|\le4$---nearly-rigid size tuples where one or two
survivors outweigh a minuscule baseline. Cascade depths re-confirmed on
the same scope: medians $2$ everywhere, maximum $4$ (zoo random-$40$).

\begin{lemma}[the free reflection fiber]\label{lem:reflect}
The map $x\mapsto-x$ on $\Z_{p^2}$ preserves every bad set
$B_w=\{x:\wn{wx}\le c\}$ and maps the fiber $F_j$ onto $F_{p-j}$. Hence,
for every unit system, the covering condition at fiber $j$ is
equivalent to the reflected covering condition at fiber $p-j$; in
particular, for every anchor solution $a\in\mathrm{Sol}_0$ with drift
$\mu$, the reflected placement covers fiber $-j_0$ with the same drift:
the fiber with dilation $\lambda\equiv-1$ has survival fraction exactly
$1$ and is never a cut.
\end{lemma}

\begin{proof}
$\wn{w(-x)}=\wn{wx}$, so $B_w=-B_w$. For $x=j+pt\in F_j$,
$-x\equiv(p-j)+p(p-1-t)\in F_{p-j}$: the involution
$t\mapsto p-1-t$ carries the footprint $B_u\cap F_j$ onto
$B_u\cap F_{p-j}$. Reflections carry arcs to arcs (an AP with canonical
difference $d$ reflects to an AP with the same canonical difference),
sizes transfer (the three-tier arc law is even in $a_0$), and the whole
fiber covering reflects to a fiber covering; the anchor's drift is a
property of the composite configuration, which reflects with it.
\end{proof}

Measured confirmation: survival fraction exactly $1$ at
$\bar\lambda=1$ in \emph{every} family at \emph{every} committed cell
($840/840$ for the zoo chain family $(1,1,1,1)$; $6/6$, $12/12$,
$24/24$ at $T=8$-critical). The resonance budget H2r is then derived
modulo the one fitted constant: the free fiber contributes the $1$, the
$\bar\lambda\le\Lambda_0$ dilations among the fiber set contribute
$\le 2\Lambda_0|U|/(p-1)\approx2\Lambda_0(T-2)/T$, and $\Lambda_0\approx4$
is where the measured cuts turn on---data-chosen, and disclosed as
such.

The case class whose structure is exactly known is the all-interval
chain family $(1,\dots,1)$---the number-one zoo family by placement
mass, with $123{,}240$ solutions at $p=17$ over all $32$ size tuples.
For it, the dilate-sparsity reduces to lattice-box counting between
junction boxes, with the load-bearing invariant the overlap-sum
identity.

\begin{lemma}[overlap-sum identity]\label{lem:overlapsum}
For the all-interval family and any size tuple $s$ with
$\Sigma s=p+\mathrm{ov}$: every covering placement, written in its
sorted cyclic order $\sigma$, has the overlap vector
$o_i=s_{\sigma(i)}-(a_{\sigma(i+1)}-a_{\sigma(i)})$ with
$\sum_i o_i=\mathrm{ov}$ exactly, $o_i\le s_{\sigma(i)}$, and every
inter-start gap $\le s_{\max}$ (the reach bound: the point before the
next start is covered by some arc starting at or before the current
one, whose reach is $\le s_{\max}$ past it).
\end{lemma}

\begin{lemma}[chain-box parameterization]\label{lem:chainbox}
$\mathrm{Sol}(s)$ is contained in the union over the $(m-1)!$ cyclic
orders $\sigma$ of the affine-unimodular images of the boxes
$\{o\in\Z^m:\ \sum o=\mathrm{ov},\ s_{\sigma(i)}-s_{\max}\le o_i\le
s_{\sigma(i)}\}$; on the $o\ge0$ stratum the parameterization is exact
up to start ties, with count
$\sum_\sigma\#\{o\ge0:\ \sum o=\mathrm{ov},\ o_i\le s_{\sigma(i)}\}$.
Consequently $|\mathrm{Sol}|\approx\mathrm{ov}^{\,m-1}$---the volume
law with exponent $m-1=T-4$ and constant derived from the overlap
budget, $\mathrm{ov}\approx(T-6)k$.
\end{lemma}

Both lemmas are machine-checked on the full solution set: $0$
violations of either bound in all $123{,}240$ solutions; the pure-chain
($o\ge0$) stratum is $120{,}960$ of them ($98.2\%$), the nested excess
$2{,}280$ ($1.8\%$); the count $131{,}520$ of (order, $o$) pairs
against $123{,}240$ placements is the tie/degenerate factor; the top
size tuple $(5,5,5,5,5)$ gives $10{,}080$ pairs against the committed
maximum $8{,}880$.

\begin{theorem}[chain-class dilate reduction]\label{thm:chaindilate}
For the all-interval family, the two-fiber survivor set at dilation
$\lambda$ is contained in
\[
\bigcup_{\sigma,\sigma'}\bigl\{o\in\mathrm{Box}_\sigma:\;
\lambda\,\tilde U_{\sigma,\sigma'}\,o+\tilde v\in
\mathrm{Box}'_{\sigma'}\pmod p\bigr\},
\]
where $\tilde U=A_{\sigma'}^{-1}A_\sigma$ is integer-unimodular
($A_\sigma$ the lower-triangular chain matrix) and the boxes have
sides $\le 2s_{\max}$. Substituting $u=\tilde Uo$ (whose box has sides
$\le 2m\,s_{\max}$), each term is bounded by
Lemma~\ref{lem:spacing} applied per coordinate:
\[
\mathrm{survivors}(\lambda)\;\le\;(m!)^2
\Bigl(1+\frac{4m\,s_{\max}+1}{b_{\min}(\lambda,\,2s_{\max})}\bigr)^{m}.
\]
Consequently the chain-class per-fiber cut is $k$-free and polynomial
in $b_{\min}$: it exceeds any fixed $\gamma$ once
$b_{\min}(\lambda,2s_{\max})\ge B(T)$, and the fibers failing this are
$\le B(4s_{\max}+1)$ in number---a vanishing fraction of $U$ as $k$
grows. Pulling the $\sigma'$-sum condition back through the dilation
gives one affine congruence
$\lambda\,(w^\top_{\sigma,\sigma'}o)\equiv c(\sigma,\sigma')$ per
order pair; for matched sizes it degenerates to the $o$-independent
congruence $\lambda\,\mathrm{ov}\equiv c$---either every anchor
candidate satisfies it or none does: the instant-kill dichotomy,
measured at $\bar\lambda\in\{2,4\}$ (survival $0$) against the partial
resonance at $\bar\lambda=3$ (survival $120/840$) in the zoo chain pair
data.
\end{theorem}

\begin{proof}
\emph{Step 1} (the parameterization): with $a_1=0$ the sorted cyclic
order $\sigma$ of the starts has $\sigma(1)=1$; the gaps
$g_t=x_{t+1}-x_t$ partition $[0,p)$, so $\sum g_t=p$ and the overlap
vector $o_t=s_{\sigma(t)}-g_t$ sums to $\mathrm{ov}$
(Lemma~\ref{lem:overlapsum}). For the reach bound: the point
$y=x_{t+1}-1$ is covered by some arc $[x_i,x_i+s_i)$ with
$x_i\le x_t$ (no start lies in $(x_t,x_{t+1})$), whence
$x_{t+1}\le x_i+s_i\le x_t+s_{\max}$. Conversely, given $(\sigma,o)$ in
the box with $\sum o=\mathrm{ov}$, the chain recursion
$x_{t+1}=x_t+s_{\sigma(t)}-o_t$ from $x_1=0$ reconstructs the starts
by a lower-triangular integer map $A_\sigma$ with determinant $\pm1$.
\emph{Step 2} (the dilate pullback): a surviving pair
$(a,a')$ with $a\in\mathrm{Sol}_0(s)$ and
$a'=\lambda a+\tau\in\mathrm{Sol}_j(s')$ writes
$a=A_\sigma o+b_\sigma$ and $a'=A_{\sigma'}o'+b_{\sigma'}$, so
$o'=\lambda\tilde U o+\tilde v$ with $\tilde U=A_{\sigma'}^{-1}A_\sigma
\in\mathrm{GL}_m(\Z)$.
\emph{Step 3} (the box count): the substitution $u=\tilde Uo$ is a
bijection of $\Z^m$ carrying the coordinate box (sides $\le s_{\max}$)
into the box of sides $\le 2m\,s_{\max}$; dropping the sum-slice
(enlarging), the count factors per coordinate and each factor obeys
Lemma~\ref{lem:spacing}: $\le|I''|/b_{\min}(\lambda,s'_{\max})+1$.
Multiplying over the $m$ coordinates and the $((m-1)!)^2$ order pairs
gives the displayed bound.
\emph{Step 4} (the sum congruence): for actual solutions
Lemma~\ref{lem:overlapsum} gives $\sum o'=\mathrm{ov}'$ as integers;
pulling through the affine expression,
$\lambda\,(w^\top o)\equiv\mathrm{ov}'-1^\top\tilde v\pmod p$ with
$w=\tilde U^\top 1$, a necessary congruence per order pair;
$o$-independent for matched sizes.
\end{proof}

What the theorem does and does not give. It gives the exact structural
account of the measured H2 phenomenology for the dominant family: the
instant kills (the sum-congruence failure), the confinement of the
resonance to small $\bar\lambda$ (the $b_{\min}$ regime), the
$k$-freeness of the cut, and the volume law that fixes the depth
formula's leading constant. It does \emph{not} give the sharp uniform
$\gamma_0=O(1)$ with clean constants---the order factors
($(m!)^2$, $4m\,s_{\max}$) are loose by orders of magnitude at $T=8$,
where the measured mid-zone cut is an instant kill and the bound gives
$\approx1$; the route to sharpness is to count \emph{on the slice}
(keep $\sum o=\mathrm{ov}$ in the lattice count---it is load-bearing;
dropping it is what costs the constants), a geometry-of-numbers lemma
in the genre of successive minima. And it does not cover the
distinct-difference families, whose junction structure is not a chain:
that is H1g, the subject of the next subsection.

\subsection{The volume law at \texorpdfstring{$T=8$}{T=8}: the census, renormalization, and the strand census}\label{sec:scale-h1}

H1g is the volume law with uniform margin,
$|\mathrm{Sol}(f,\mathrm{sz})|\le c_1\,\mathrm{ov}^{\,m-1}$ for the
$\pm$-distinct families, with $c_1\le\bigl(T/(T-6)\bigr)^{m-1}e^{-\eta}$
for an $\eta>0$ uniform in $(T,k,f,\mathrm{sz})$---the margin clause
that keeps Proposition~\ref{prop:Hq}'s denominator positive. The
bare-$c_1$ form below is the measured content; the margin clause is
the open content. It was attacked at the smallest rung where the
pattern is visible, $T=8$, by full census where the cell admits it:
$p=17$ ($k=2$, $\epsilon=1$; all $840$ families, all $32$ size tuples,
$\mathrm{ov}\in[3,8]$; per-family totals median $2{,}981$, $90$th
percentile $9{,}633$; per-(family, size) maximum $876$ at
$(5,5,5,5,5)$, $\mathrm{ov}=8$); $p=19$ ($k=2$, $\epsilon=3$; all
$1{,}680$ families; admissible $960$, exactly the pool saturation of
Table~\ref{tab:pool} reproduced by an independent code path; totals
median $102$, $90$th percentile $645$, maximum $801$; per-size maximum
$108$ at $\mathrm{ov}=6$); and $p=29$ ($k=3$, $\epsilon=5$;
\emph{targeted}: every cascade-containing set, the doubling-path
quadruples, the anchors, and a $120$-set random sample---$138$ of the
$715$ difference sets, $710$ families of the $17{,}160$ total, a scope
disclosed and not claimed beyond; $20/138$ sets admissible, consistent
with the full-scope $0.133$; per-size maximum $186$ at
$(8,8,8,8,8)$, $\mathrm{ov}=11$). Every decision procedure was
validated against an independent direct brute force (plain $p^4$
enumeration, no cascade, no pruning) before any number was cited:
$9/9$ keys match exactly across the three primes, including the chain
controls $8{,}880/4{,}680/29{,}760$ and zero keys.

\begin{lemma}[renormalization invariance]\label{lem:renorm}
Let $D=\{1\}\cup S$ be the difference multiset of a census family ($S$
a $\pm$-distinct $(m-1)$-set), and for $\delta\in D$ let
$\mathrm{renorm}_\delta(S)=\{\mathrm{canon}(d\,\delta^{-1}\bmod p):
d\in D,\ d\ne\delta\}$, $\mathrm{canon}(x)=\min(x,p-x)$. Then the total
covering count summed over all size assignments is equal for $S$ and
$\mathrm{renorm}_\delta(S)$:
$\sum_{\mathrm{sz}}|\mathrm{Sol}(S,\mathrm{sz})|
=\sum_{\mathrm{sz}}|\mathrm{Sol}(\mathrm{renorm}_\delta(S),
\mathrm{sz})|$.
\end{lemma}

\begin{proof}
The covering condition is invariant under the global dilation
$\varphi(x)=\delta^{-1}x$ of $\Z_p$ (a bijection carrying an arc
$a+d\,[0,s)$ onto the arc $\delta^{-1}a+(d\,\delta^{-1})[0,s)$). Apply
$\varphi$ to a covering placement of the family $S$ with size
assignment $(s_d)_{d\in D}$: the image is a covering placement of the
family with difference multiset $\delta^{-1}D$, in which the arc that
had difference $\delta$ now has difference $1$---it is the normalizer
of the transformed system. Re-anchor at that arc's start (covering is
translation-invariant) and canonicalize each difference
$d\mapsto\mathrm{canon}(d)$ (an arc with difference $-d'$ is the
reflection, $x\mapsto-x$ composed with a translation, of an arc with
difference $d'$; reflections preserve coverings). The result is a
census-normal-form covering placement of
$\mathrm{renorm}_\delta(S)$ with the sizes carried along. Summing over
all size tuples removes the bookkeeping dependence on which arc
received which size (the admissible size-tuple set is
permutation-closed), and the map is a bijection at each size
assignment.
\end{proof}

Verified exactly: at $p=17$ the $35$ difference sets partition into
$7$ renormalization orbits (each of size $5$, one per choice of
$\delta\in D$), and the $7$ volume classes are precisely these orbits;
at $p=19$ the $70$ sets form $14$ orbits ($8$ positive, $6$ zero), and
the $9$ volume classes are exactly the unions of orbits. The volume is
therefore a function of the dilation class of the difference multiset
alone---not of the family key. This cuts the family space $5$-to-$1$
in general and identifies the correct equivalence for any future
classification: at $T=8$ the census object is not the
$\binom{D-2}{m-1}$ sets but their renormalization orbits. It is a
dilation-invariance statement in the spirit of the class reduction at
prime moduli (Lemma~\ref{lem:class}) and the census law of
Conjecture~\ref{conj:SL}: it identifies the equivalence; it does not
count. For the cascade the consequence is load-bearing: the two-fiber
cut at dilation $\lambda$ compares $\mathrm{Sol}_0$ with a dilate of
$\mathrm{Sol}_j$, and Lemma~\ref{lem:renorm} says the volumes involved
are dilation-class invariants---the cut law's input volumes are orbit
functions, and H1g needs constants uniform over orbits, not over keys.

\begin{table}[htbp]
\centering
\small
\caption{The volume law at $T=8$: the worst $\pm$-distinct
renormalization orbit against the chain benchmark, at the top size
tuple. Normalized: the maximum divided by $\mathrm{ov}^4$. The uniform
bound $|\mathrm{Sol}|\le\mathrm{ov}^4$ (that is, $c_1=1$) holds at
every census prime with margin $4.7\times/12\times/79\times$,
\emph{growing in $k$}; the distinct class sits
$10\times/43\times/160\times$ below the chain class at equal
$\mathrm{ov}$. The $p=29$ top orbit is identified within the targeted
scope; a family outside it would have to exceed $\mathrm{ov}^4$
outright to change the law's direction, and the admissibility collapse
argues against it.}
\label{tab:volume}
\begin{tabular*}{\textwidth}{@{\extracolsep{\fill}} r l r r r r}
\toprule
$p$ & top distinct orbit & max @ $\mathrm{ov}$ & normalized &
chain @ tuple & ratio \\
\midrule
$17$ & $\{1,2,3,4,8\}$ & $876$ @ $8$ & $0.214$ & $8{,}880$ & $0.099$ \\
$19$ & $\{1,2,4,5,8\}$ & $108$ @ $6$ & $0.083$ & $4{,}680$ & $0.023$ \\
$29$ & $\{1,2,4,7,8\}$ & $186$ @ $11$ & $0.013$ & $29{,}760$ & $0.0063$ \\
\bottomrule
\end{tabular*}
\end{table}

The reading of Table~\ref{tab:volume} is fourfold. \emph{(i)} The
uniform bound with $c_1=1$ holds at every measured prime, with margin
growing in $k$ (and tightening with the $\epsilon$ window at fixed
$k$). \emph{(ii)} The distinct class is uniformly \emph{thinner} than
the chain class at equal $\mathrm{ov}$, by a factor growing from
$10\times$ to $160\times$. \emph{(iii)} The direction of risk for H1g
is therefore not a fat counterexample family but proof coverage---the
opposite of an obstruction. \emph{(iv)} The honest caveats: three
primes, two $k$-values, both $\epsilon$ classes at $k=2$ but only one
at $k=3$; the $p=29$ top orbit rests on the targeted scope; the
measurements fix the law's direction and margin, not its proof.

\begin{proposition}[the threshold form of H1g]\label{prop:threshold}
The per-coordinate bound $|\mathrm{Sol}(f,\mathrm{sz})|\le p^{\,m-1}$
is trivial at every cell of the measured family (the
moving-start space is $\Z_p^{\,m-1}$), and $p/\mathrm{ov}\to T/(T-6)$
because $\mathrm{ov}=(T-6)k+O(1)$ against $p=Tk+\epsilon$, so the
trivial bound realizes
\[
c_1^{\mathrm{triv}}=\bigl(T/(T-6)+o(1)\bigr)^{\,m-1}
\]
in the normalization of the volume law---exactly the value at which
the depth denominator $(m-1)\log\bigl(T/(T-6)\bigr)-\log c_1$ of
Theorem~\ref{thm:SC} vanishes. At $m-1=T-4$ the threshold is $256$,
$243$, $244.1$ at $T=8/9/10$: near-constant across rungs, so neither
the cascade's task nor the volume law's softens at higher $T$.
\end{proposition}

\begin{proof}
The first inequality is the per-coordinate count: a solution is a
choice of the $m-1$ moving starts, each in $\Z_p$. The second is the
size window of Lemma~\ref{lem:rungT}: the $m=T-3$ arcs of the
two-value window give $\mathrm{ov}=\Sigma s-p=(T-6)k+O(1)$ against
$p=Tk+\epsilon$. Dividing,
$|\mathrm{Sol}|\le(p/\mathrm{ov})^{\,m-1}\mathrm{ov}^{\,m-1}$, which is
the volume law with $c_1=(p/\mathrm{ov})^{\,m-1}$; substituting this
$c_1$ in the depth formula of Theorem~\ref{thm:SC} zeroes its
denominator, so the trivial bound purchases no cascade depth at all.
\end{proof}

The threshold form restates what H1g must purchase: not the free
bound $|\mathrm{Sol}|\le p^{\,m-1}$ but the \emph{margin} over it,
growing in $k$. Measured, the margins below the trivial box are
$95\times/1{,}207\times/3{,}803\times$ at $p=17/19/29$ (top distinct
family at its top size tuple: $17^{4}/876$, $19^{4}/108$,
$29^{4}/186$---the maxima of Table~\ref{tab:volume}); the committed
margins over $\mathrm{ov}^{4}$ ($4.7\times/12\times/79\times$) are
the same fact in volume-law units. The form also unifies the program's
constant-quality obligations under one numerical criterion: every
volume constant must beat $\bigl(T/(T-6)\bigr)^{\,T-4}\approx250$.
Every measured constant clears it---the chain class's
$2.0$--$3.6$ by roughly two orders of magnitude, the distinct class by
the factors above---and no proved constant yet does (the chain
constants of Theorem~\ref{thm:chaindilate} are loose by orders of
magnitude). The sharp-constants question of \S\ref{sec:scale-h2} and
H1g are one criterion stated twice.

The mechanism of the top class was probed on the full solution sets.
Every solution needs all five arcs (essential-arc pattern
$(1,1,1,1,1)$ in $876/876$): no spectator redundancy. The
overlap-budget identity
$\sum_{i<j}O_{ij}=\mathrm{ov}+\sum_x\binom{o(x)}{2}$ holds exactly
($0$ violations), with pairwise overlaps small (mean $0.5$--$1.5$,
maximum $2$--$3$) and under the spacing caps with room: no pair
carries the covering; the arrangement is collective. Coordinate
projections of the solution set at the probed families sit far from
full ($[16,17,14,12]$ for the top family; $[9,9,7,8]$ for the minimum
class)---per-family statistics that initially read as a $T=7$-style
pinning law persisting at the zoo cell; the exhaustive key-level check
of the character-system paragraph below shows that reading is an
artifact, not a law. And the block geometry exhibits
the mechanism: the interval covers $[0,s_1)$; the $d=2$ arc covers
every other point of a window of span $2s-1$; the $d=4$ arc every
fourth point of a span $4(s-1)$; the $d=8$ arc is the antipodal
two/three-block sampler; the circle closes by interleaving these
against the interval and each other. The top class is the
\emph{binary cascade} $\{1,2,4,8\}$ plus a fifth difference---at
$p=17$ the orbit swallows every fifth value $\{3,5,6,7\}$; at $p=19$
the fifth is $5$; at $p=29$ it is $7$ (the observed pattern
``fifth $=\epsilon+2$'' holds at the three primes measured; three
points, flagged as such, not a law). At $p=29$, where $2$ is a
primitive root, the canonicalized doubling orbit visits every
difference and only cascade-containing quadruples survive among the
path quadruples: resonance structure is necessary for admissibility at
$k=3$, consistent with the pool collapse.

\paragraph{The character system: the corrected coordinates.}
The projection reading above was tested as a law and is false in both
naive forms. Call the \emph{character marginals} of a solution set the
projection sizes along the coordinate functionals of the moving-start
space. The two natural invariance readings---(V$'$) the multiset of the
four singleton marginals is a renormalization-orbit invariant, and
(V$''$) the same for the six pairwise-difference marginals---are both
false, each caught by machine check (V$'$: profile multisets differ
across orbit members at both full primes; V$''$: $0/32$
$\sigma$-matched size tuples agree at $p=19$). The mechanism was
pinned by an explicit construction of the renormalization bijection
(at $p=19$, $S=(2,4,5,8)$, $\delta=5$, with image $S'=(3,4,6,8)$):
the image of each moving slot is $\delta^{-1}$ times a source slot
coordinate \emph{minus the anchor coordinate} $a_\delta$, plus a
size-dependent constant, and the new normalizer is
$\delta^{-1}a_\delta$---every moving slot is coupled to the anchor.
Consequently the image's slot singletons are the source's anchor pairs,
and the image's pairwise differences (in which the anchor cancels) are
the source's non-anchor singletons and pairs: neither subsystem is
closed, and the union is.

\begin{lemma}[the extended character system]\label{lem:char10}
Let $X=\{e_i\}\cup\{e_i-e_j:\ i<j\}$ on the moving-start space
$\Z_p^{\,m-1}$, so $|X|=(m-1)+\binom{m-1}{2}$ (ten functionals at
$T=8$). Under the renormalization bijection of Lemma~\ref{lem:renorm},
each functional of $X$ maps to a nonzero scalar multiple of a
functional of $X$ plus a size-dependent constant. Consequently the
multiset of marginal sizes
$\{|\pi_\chi(\mathrm{Sol}(S,\mathrm{sz}))|:\chi\in X\}$, aggregated
over the permutation-closed set of size tuples, is a
renormalization-orbit invariant.
\end{lemma}

\begin{proof}
The closure statement is the anchor-mixing computation just recorded,
verified symbolically and then end-to-end at $p=19$ for
$S=(2,4,5,8)$: all $801$ images of $\mathrm{Sol}(S,\mathrm{sz})$ land
in $\mathrm{Sol}(S',\mathrm{sz}\circ\sigma)$, injectively, at every
size tuple, with per-key counts matching under the explicit $\sigma$
($0$ mismatches)---which also confirms the per-assignment-bijection
step of Lemma~\ref{lem:renorm}'s proof at a stronger level than the
committed record (totals were verified; assignments now are). The
invariance follows: each character's marginal size is carried to
another character's marginal size up to the size permutation, so the
aggregate multiset is constant along the orbit. Verified on the full
top orbits at both primes---$p=17$: the orbit of $(2,3,4,8)$, five
members, $9{,}633$ solutions per set; $p=19$: the orbit of
$(2,4,5,8)$, five members, $801$ per set---with the aggregate
ten-functional multisets identical across members at both primes; and
at $k=3$, the two admissible $p=29$ probe sets $(2,4,7,8)$ and
$(3,4,8,13)$ form one orbit ($\delta=7$) with equal totals $1{,}311$.
\end{proof}

Two consequences. First, the coordinate system for the census is now
fixed: any class-level (orbit-level) law must be stated on
$X$---not on absolute starts, not on pairwise phases alone---so the
symmetry analysis of the census object is complete
(Lemma~\ref{lem:renorm} fixes the equivalence, and
Lemma~\ref{lem:char10} the coordinates in which a law can be stated).
Second, the \emph{marginal route into the strand census is measured
dead} at the zoo cell. Across $1{,}678$ admissible keys (exhaustive at
$p=17$ and $p=19$, targeted at $p=29$; every key inside the committed
census scope matches it, $1{,}654/1{,}654$, $0$ mismatches), the
projection envelope over all keys and all ten characters is $p$ itself
at both $k=2$ primes and $27/29$ at $k=3$ (fullness fractions
$1.000/1.000/0.931$); a sufficient marginal law on $X$---
$|\pi_\chi|\le\gamma_0\,\mathrm{ov}$ with
$\gamma_0^{\,m-1}<\bigl(T/(T-6)\bigr)^{\,m-1}$---would imply the
volume law, since $\mathrm{Sol}$ is contained in the product of its
singleton marginals, but it requires the fullness fraction to fall
below $\gamma_0(T-6)/T$ asymptotically (below $75\%$ already at
$\gamma_0=3$), and no downward trend is in evidence at the three
committed primes. The top-family profiles are non-monotone in
$(k,\epsilon)$ ($[12,17,17,16]\to[4,4,9,8]\to[9,7,10,27]$; the free
character $27$ at $p=29$ is the $d=8$ antipodal sampler of the block
geometry above): family- and window-dependent statistics, not a law.
The only invariant sparsity is \emph{joint}: the median density of
$\mathrm{Sol}$ inside its own marginal box is $1.5\%/7.4\%/6.9\%$ at
the three primes---the pointwise no-holes condition along $I$, which
is what the strand language below encodes.

\paragraph{The strand census: where H1g now lives.}
The reduction runs in three steps, each elementary. \emph{Step 1}
(the interval spine): the complement of the normalizer interval is the
interval $I=[s_1,p)$ of length $\ell=p-s_1\approx(T-2)k$; every
covering of $\Z_p$ by the interval plus the $m-1$ moving APs restricts
to a covering of $I$ by the $I$-restrictions of the moving APs (plus
the two boundary points at the wrap), and the restriction map is
injective up to the boundary bookkeeping, so an upper bound on the
$I$-restricted count bounds $|\mathrm{Sol}|$. The moving arcs' total
size is $(m-1)\cdot2k\approx(T-3)\cdot2k$ against
$\ell\approx(T-2)k$: a near-exact cover with slack $\mathrm{ov}$.
\emph{Step 2} (the strands): the $I$-restriction of an AP with step
$d$ is a union of \emph{strands}---maximal runs of consecutive $t$
whose points $a+td$ stay inside $I$; within a strand the points sit $d$
apart. Covering consecutive points of $I$ requires strands of
different steps to interleave; a pair $(d,d')$ interleaves only at
compatible phases, and the phase-compatibility condition per pair is
a three-distance-type statement about $\lambda=d'/d$---exactly the
Lemma~\ref{lem:E} genre, on the pair lines $td-ud'\equiv\delta$ that
already govern the pairwise overlaps. \emph{Step 3} (the arrangement
form): every covering induces a cyclic \emph{arrangement}---the
maximal runs of $I$ and, for each run, the subset of strands carrying
it, with the handoffs between runs subject to the pair-phase
conditions---and the volume decomposes as
$|\mathrm{Sol}|\le\sum_{\text{arrangements}}(\text{phase freedom of
the arrangement})$, with boundary corrections.

What is missing for a proof is the uniform count: bound the number of
arrangements (the handoff patterns) and the phase freedom per
arrangement, uniformly over the resonance classes of the difference
set---equivalently, by Lemma~\ref{lem:renorm}, over the
renormalization orbits. The chain class is the degenerate resonance
class where all steps coincide and the strand picture collapses to the
telescoping of Lemmas~\ref{lem:overlapsum}--\ref{lem:chainbox}; the
binary cascade is the pure ratio-$2$ class; a general orbit mixes
ratio classes. The three families of necessary conditions already
available---the overlap budget, the pair spacing caps, and the
last-arc fit count---do \emph{not} discriminate: random placements
typically satisfy the budget (the expected total pairwise overlap
$\approx12$ exceeds $\mathrm{ov}$ at $p=17$). The discriminating
structure is the pointwise no-holes condition along $I$, which is what
the strand language encodes. We record the reduction honestly: it is a
reduction to a named, bounded combinatorial object, the \emph{strand
census}; the injectivity of the restriction map and the boundary
corrections are machine-consistent at the census cells but are not
proved in general, and the strand census is therefore a
\emph{plausible} reduction of H1g, not an established equivalence. The
work is class-by-class (per resonance type), and uniformity across
classes is itself a lemma---the precise shape of the remaining distance.

\paragraph{The transfer-matrix session: the third route, closed by
theorem.}
The third candidate route---a transfer-matrix or automaton model of the
joint density under the arc-covering constraint---was executed as a
proof attempt under a pre-committed success criterion: a bound on the
whole-cascade survival fraction, survivors $\le\rho^{\mathrm{depth}}
\cdot$ initial mass, with $\rho\le 1-\epsilon$ for some $\epsilon>0$
independent of $k$. Two specifications were fixed before the session,
as the route's design demanded. \emph{The state space}: the covering
constraint at fiber $j$ is membership
$c(j)\in\mathrm{Sol}(\mathrm{sz}(j))$---a function of the placement
state alone; the cascade is non-consuming (the alive-set semantics is
a pure intersection), so the constraint is Markovian in the natural
state $(\Z_p)^{m-1}$, and the theorem below is invariant under state
augmentation in any case, so no resolution of this specification could
change the verdict. \emph{The comparison target}: the whole-cascade
fraction, with both factors' $k$-dependence acknowledged---the initial
mass is $|\mathrm{Sol}_0|$, and the per-step ratios inherit the
$k$-scaling of the forbidden sets.

\begin{lemma}[the monomial dichotomy of the cascade]\label{lem:monomial}
Fix the rung, the prime, and any sequence of fibers $j_1,\dots,j_D$
from the uncovered set with solution sets
$\mathrm{Sol}_t=\mathrm{Sol}(\mathrm{sz}(j_t))$, and let $\sigma_t$
be the affine bijection between consecutive fiber placements of
Lemma~\ref{lem:twofiber}. Then: \emph{(i)} every constrained
transition $P_{\mathrm{Sol}_{t+1}}\,\sigma_t\,P_{\mathrm{Sol}_t}$ is
a partial permutation matrix (at most one nonzero entry, equal to $1$,
per row and per column), and so is every product along the chain; its
spectrum is $\{0\}$ (the directed paths) together with the roots of
unity of the surviving cycle lengths. \emph{(ii)} The affine maps
compose to the identity around any closed fiber loop (each re-expresses
the fixed drift $\mu$; the dilation ratios telescope), so the cyclic
transfer product is a diagonal $0/1$ matrix whose trace is the number
of placements surviving the loop---the cascade's own survivor
count---and whose spectral radius is $1$ if that count is positive and
$0$ otherwise. \emph{(iii)} Consequently every constrained transfer
matrix obtainable in the framework---open chain, cyclic, homogenized
single pair, or state-augmented---has spectral radius exactly $0$ or
$1$: never an interior value, at any $k$.
\end{lemma}

\begin{proof}
\emph{(i)} $\sigma_t$ is a bijection; composing with coordinate
projections leaves at most one nonzero per row and per column, and
partial permutations are closed under products. The digraph of a
partial permutation has in- and out-degree at most $1$, hence
decomposes into directed paths and cycles; paths give nilpotent
blocks, and a cycle of length $\ell$ gives a cyclic permutation block
whose spectrum is the $\ell$-th roots of unity.
\emph{(ii)} Each $\sigma$ is the change of drift coordinates
$c(j')=(j'/j)\,c(j)+c^0(j')-(j'/j)\,c^0(j)$
(Lemma~\ref{lem:twofiber}); composing around a loop telescopes the
dilation ratios to $1$ and the translations to $0$, and deletion only
zeroes diagonal entries. \emph{(iii)} Augmentation preserves
determinism (the register update is a function of the state), so the
augmented constrained matrix is again a partial permutation, and
parts (i)--(ii) apply verbatim.
\end{proof}

The machine record runs at three levels. \emph{Generic}:
$1{,}500/1{,}500$ random constrained systems (random bijections with
random deletion sets) satisfy the structural dichotomy---the
partial-permutation premise, and acyclic if and only if
nilpotent---and $12{,}000$ exact trace fingerprints
$\mathrm{tr}(Q^t)=\sum_{\ell\mid t}\ell\,c_\ell$ (over $300$ systems,
$t\le 40$, integer-exact) match the cycle inventory with zero
mismatches. \emph{Committed objects}: six committed families (the zoo
chain and top-distinct orbits at $p=17$ and $p=19$; the two committed
admissible $p=29$ sets), $66$ homogenized pair systems---the route's
own model, iterating a single pair map $x\mapsto\lambda x+\tau$ under
the deletion set $\mathrm{Sol}_j$---of which exactly the $6$
reflection pairs ($\bar\lambda=1$) have $\rho=1$ and all $60$ others
have $\rho=0$; the loop-composition identity of part (ii) holds
algebraically and numerically ($500$ random points per family); the
automaton's per-step survivor counts equal the committed $\mu$-space
cascade convention at every step (the chain family's depth $2$
reproduces the committed zoo record); and the end-to-end ground
truth---$[\sigma(x)\in\mathrm{Sol}_j]$ if and only if the direct
fiber-covering test accepts $\sigma(x)$---holds on $2{,}800/2{,}800$
sampled checks, with the committed volume margins reproduced exactly
(top $|\mathrm{Sol}|=876/108/186$, margins
$95\times/1{,}207\times/3{,}803\times$ over the trivial box).
\emph{The reflection fiber in automaton language}: the two-fiber
survival fraction at $\lambda=-1$ is exactly $1$ at all six families
across all three primes, and the homogenized reflection sub-shift
closes the whole of $\mathrm{Sol}_j$ into $2$-cycles---survival flat
forever: Lemma~\ref{lem:reflect} as a statement about the instrument
itself, the deletion at the free fiber being empty.

What the instrument measures at the committed cells is equally sharp.
The iterated-pair survival counts collapse as $[840,120,0]$,
$[224,1,0]$, $[4{,}680,240,0]$---piecewise extinction within three
iterates, the intermediate values set by maximal lattice runs of
solution points along the drift orbits (runs of length at most $2$
everywhere; extinction confirmed at every horizon to $12$)---the
three-distance genre of Lemma~\ref{lem:E}, not geometric decay with
an interior rate. And the $k$-scaling that the obligations need (the
margins $95\times/1{,}207\times/3{,}803\times$) lives in
$|\mathrm{Sol}|/p^{m-1}$, in the census counts, and in no spectral
quantity: the spectral outputs are binary at every prime, with no
$k$-trend because there is nothing to trend.

Nor does any reformulation of the instrument escape the dichotomy. A
coarse-grained or weighted automaton (states as classes, class
multiplicities as weights) can carry an interior spectral radius
numerically, but for non-negative matrices the radius dominates every
directed cycle's geometric mean, so $\rho\le 1-\epsilon$ would force
uniform decay of the weight products along drift orbits---and the
weights are census counts, so this is the margin itself entering as an
input rather than an output; the converse fails besides (the $2\times2$
witness $\left(\begin{smallmatrix}0&0.9\\0.9&0.4\end{smallmatrix}
\right)$: every directed cycle has geometric mean $\le 0.9$ yet
$\rho\approx 1.122$---branching resurrects mass), so a second external
input, a branching control, would be needed as well. Averaged or
randomized operators collapse to the independent-extension baseline,
whose deficit is the $\gamma$-correlation---the cascade model's own
deepest unproved input. Every door leads through an external
uniform-in-$k$ estimate of exactly the margin type; none generates
one. In automaton language the whole-cascade survival fraction is the
product of the consecutive conditional ratios $r_t=|S_t|/|S_{t-1}|$,
each bounded by $\gamma_0\,|\mathrm{Sol}_{j_t}|/p^{m-1}$ under H2m
with the thinness of $|\mathrm{Sol}|$ supplied by H1g: under its own
model the route's output is Proposition~\ref{prop:depth}'s recursion,
and the uniform deficit it owes is the hypothesis stack it was
commissioned to prove. The instrument is the census in different
notation.

The verdict against the pre-committed trichotomy of outcomes: the
success criterion ($\rho\le 1-\epsilon$) is unattainable by
theorem---there is no interior spectral radius to bound, in any state
formulation, at any $k$; and the failure mode is not $\rho\to 1$ as
$k\to\infty$ but $\rho\in\{0,1\}$ as an identity of the lossless
dynamics. The type-mismatch reading is thereby upgraded from
observation to mechanism: bijections preserve measure exactly (the
monomial structure), deletion filters without dissipating (partial
permutations carry spectrum in $\{0\}$ together with the unit circle),
and the loop composition is the identity (the cyclic instrument
measures alive-or-dead, not decay). Margins require dissipation;
losslessness forbids generated dissipation. The reformulation is
bounded by its own fidelity. Lemma~\ref{lem:monomial} is the program's
third $k$-invariance theorem---after the spacing lemma's
$k$-cancellation and the character system's closed coordinates---and,
like them, a closure-type statement: the toolkit's signature holds
even for the instrument built to break it.

\subsection{The verification record of this section}\label{sec:scale-sound}

Every number above is either an exhaustive census with two independent
decision procedures, or a sample with its size and interval, or a
proved statement with its machine check; the records, in the order of
appearance: the pool censuses (the pruned-cascade and the
flattened-popcount procedures agree exactly at $p=17/19/23$; an
independent pure-enumeration brute force agrees on $21/21$ sampled
tuples across $p=17/19/23/31$, covering both admissible and
inadmissible branches; the $T=9$ stride samples carry full decisions,
and the $m=6$ code path shares the cascade code brute-force-verified
at four $m=5$ primes, with the external consistency check the
agreement of the projected pools with the independently measured pool
scale); the zone-resolved $\gamma$ table ($764$ pairs, five committed
cells, the identical seed reproducing the committed family
selections); the chain-class machine checks ($123{,}240$ solutions,
$0$ violations of the sum identity or the reach bound; the count
reconciliations $131{,}520$ against $123{,}240$ and $10{,}080$ against
$8{,}880$); the reflection fiber (survival exactly $1$ at every
$\bar\lambda=1$ pair, every family, every cell); the
renormalization-orbit law (machine-checked set-by-set at $p=17$ and
$p=19$); the volume census and its ground truth ($9/9$ brute-force
keys; the admissibility reproductions $840/840$ and $960/1{,}680$; the
order-independence of volume across all $24$ orderings); the
essential-arc, budget-identity, and cap checks on the probed families;
and the cascade depths at the committed cells (medians $2$, maximum
$4$); and the character-system record of \S\ref{sec:scale-h1} (the
rigidity census over $1{,}678$ admissible keys with every key inside
the committed census scope matching it exactly, $1{,}654/1{,}654$, $0$
mismatches; the per-assignment renormalization bijection, $801/801$
images injective with per-key counts matching under the explicit
$\sigma$; the falsifications of V$'$ and V$''$; the character-system
aggregate invariance on both full top orbits; the chain controls
$8{,}880/4{,}680/29{,}760$ reproduced by the same code path). The
transfer-matrix session's record: the generic dichotomy and trace
fingerprints ($1{,}500$ systems, $12{,}000$ comparisons, zero
mismatches); the $66$ homogenized pair systems on the six committed
families with the closed-form orbit inventories verified by full
traversal and the finite-horizon survival counts verified by direct
simulation (zero mismatches); the loop-composition identity
(algebraic and numeric, $500$ points per family); the per-step
equality of the automaton's survivor counts with the committed
$\mu$-space cascade convention (zero mismatches, all steps, all
families; the committed chain-family depth reproduced); the
end-to-end fiber-covering ground truth ($2{,}800/2{,}800$); the
reflection pairs (survival fraction exactly $1$, all six families,
all three primes); and the reproduction of the committed volume
margins ($876/108/186$ against $95\times/1{,}207\times/3{,}803\times$).
The three-rung depth record of Table~\ref{tab:depthcal} is
calibration from the earlier cascade runs and is labeled as such
wherever it is used.

\subsection{The epistemic status}\label{sec:scale-status}

The program is either one proof away---H1g and the sharp H2m, both
reduced to concrete named objects, the strand census and the pair-line
count---or one systematic obstruction away: a family class whose
junction structure escapes the $\mathrm{ov}^{\,m-1}$ volume law. For
the obstruction direction the $T=8$ evidence points the other way: the
measured margin \emph{grows} with $k$, and no distinct family reaches
even a tenth of the chain-class volume at equal $\mathrm{ov}$ beyond
$k=2$. The distinction will be settled by whether the strand census
closes at $T=8$ with uniform constants, not by further measurement.

One more disclosure belongs in the record. The sequence of named
obligations has moved: the input structure, then the exhaustiveness of
the census classification, then the scaling-closure lemma itself, then
H1-general, now the strand census. Each closure has been real, and
each new obligation has been at the same scale as the previous one.
Whether this is a fractal obstruction---each gap revealing a new gap of
the same size, indefinitely---or a sequence of genuinely decreasing
obligations that will bottom out is not decidable from the present
data; after five iterations of the pattern the defensible prior is that
the reformulation has a stubborn residual, and Theorem~\ref{thm:SC} is
stated so that the residual is exactly what the hypotheses name. What
would change the verdict: a family at some $(T=8,k)$ whose per-size
volume exceeds $\mathrm{ov}^{\,m-1}$---the measured pool collapse
makes this increasingly unlikely, as admissibility itself dies at
$13\%$ by $p=29$---or a strand-census proof at $T=8$ with uniform
constants, which would make H1g a theorem at the rung and supply the
template for general $T$. No claim is made here about the number of
sessions, the amount of work, or the likelihood of either outcome; the
failure modes are stated where the obligations are.

The record since that disclosure adds a sixth attack turn and a change
in the pattern's shape. The strand census has now been attacked twice
and has survived both attacks: the necessary-condition families of
this section (the overlap budget, the spacing caps, the last-arc fit
count---non-discriminating), and the marginal-rigidity route of
\S\ref{sec:scale-h1} (the projection envelope measured full, the two
naive coordinate readings false, the corrected character system closed
but silent on the count). Both closed as route-closures, and---for the
first time in the sequence's six attack turns---the named obligation
did not move: no successor lemma was minted; what moved was the
obligation's coordinate system (Lemma~\ref{lem:char10}) and its
numerical form (Proposition~\ref{prop:threshold}). The honest reading
has tightened accordingly: the residual is not a bookkeeping artifact
awaiting the right census format; it is resistant to the toolkit that
produced every proved result in this paper. A structural fact in the
ledger explains why that may be: the proved results are closures,
invariances, exhaustiveness statements, and machine verifications,
while every remaining obligation---H1g's margin over the trivial
bound, the volume law's unproved growth in $k$, the chain class's gap
from measured to proved constants---is a growth-in-$k$ margin over a
trivial or measured baseline, and no turn of this program has proved a
statement of that type. Both closed routes, moreover, were
census-shaped attacks on what Proposition~\ref{prop:threshold} shows
is an inequality-shaped obstruction. Three candidate third routes are
on record, none attempted here and none with a stated prior:
pairwise-difference marginals in the corrected $X$ coordinates; a
combinatorial characterization of the joint support, given its
measured sparsity ($1.5\%$--$7.4\%$ median density); and a
transfer-matrix or automaton model of the joint density under the
arc-covering constraint.

The type-mismatch reading of the preceding paragraph has since been
verified statement-by-statement against the paper's full proof
inventory---every lemma, proposition, theorem, and machine record---and
it holds, with the three statements that could be read as exceptions
resolved explicitly. First, the spacing lemma (Lemma~\ref{lem:spacing})
is a packing bound whose slack over its own trivial baseline is the
factor $b_{\min}$; at the cascade's operating point ($L=Rk$ against
$p=Tk$) the bulk-fiber value $b_{\min}\gtrsim p/(2L+1)\approx T/(2R)$ is
$k$-free, so the lemma is invariant under exactly the scaling in which
the obligations' margins grow, and the ``$k$-cancellation''
(Proposition~\ref{prop:kcanc}) is a $k$-invariance of the cut,
conditional on the volume law---the cancellation of the dependence the
obligations need to grow, not a proof that it grows. Second, part (d) of
the comparison (Proposition~\ref{prop:depth}) does state a margin
growing linearly in $k$, but as a consequence of its hypothesis stack
(H1g included), not as a proved statement about the measured world.
Third, the chain-class volume law (Lemma~\ref{lem:chainbox}) computes a
count that grows in $k$, but its margin over the trivial box is the
$k$-free threshold value of Proposition~\ref{prop:threshold}; the
growth-in-$k$ content of H1g is precisely the distinct class's extra
thinness below the chain benchmark ($10\times/43\times/160\times$ at
equal $\mathrm{ov}$, measured), which the chain-class theorem does not
touch, while the proved chain constants of Theorem~\ref{thm:chaindilate}
carry order factors ($4m\,s_{\max}$ with $s_{\max}\approx2k$) whose
slack grows in $k$---the measured-to-proved constant gap already listed
among the obligations. With the reading verified, the three candidate
third routes of the preceding paragraph are typed by output:
pairwise-difference marginals in the corrected $X$ coordinates and the
joint-support characterization are census-shaped (classification
outputs; the same type as the two routes already closed), while the
transfer-matrix or automaton model is the one margin-shaped
candidate---the natural output of a transfer-operator argument is a
spectral quantity uniform in $k$, the missing statement type. The typing
assigns output types, not success probabilities, and rests on the
tool-reach half of the reading rather than a theorem; it is the reason
the automaton route is the only one the record prioritizes, and the
attempt itself is its test.

The attempt has since been run, under the discipline of a proof
attempt: the two design specifications fixed beforehand, the success
criterion pre-committed, and the verdict taken against it. The route
closed by theorem (Lemma~\ref{lem:monomial} and the analysis following
it): every spectral quantity the automaton framework can produce is
exactly $0$ or $1$, the cyclic product's trace is the cascade's own
survivor count, and the survival-fraction deficit the criterion
demanded translates exactly into the pair bound H2m and the volume law
H1g---the instrument consumes the hypotheses it was commissioned to
prove. The typing's prediction was therefore right about the output
type and wrong about its shape: the route's natural output is not a
margin in disguise but a closure about itself, and the expected
spectral gap was an artifact of the classical analogy (sub-shifts of
finite type earn their gaps through branching that a lossless,
sufficient-statistic dynamics cannot have). Three routes into the
strand census are now closed---the necessary-condition families, the
marginal-rigidity route, and the automaton---the first two by
measurement, the third by theorem. The empirical half of the reading
(S1) is verified a second time by a second instrument: the one
margin-shaped candidate, fully specified and executed, emitted exact
census computations and a $k$-invariance, and no margin. The
tool-reach half (S2) has its first proved instance---every automaton
built on a bijective drift is census-shaped---which is a mechanism,
not yet the general claim; instruments that do not inherit the drift
remain outside its scope, and nothing here is claimed about them. With
this, the program's state is fixed, and it is the state the
conditional theorem was stated to name: Theorem~\ref{thm:SC} with its
ledger; the type-mismatch diagnosis, now with its mechanism
(losslessness forbids generated decay); and a named residual---the
growth-in-$k$ margin, uniform over families, size tuples, and
$k$---that the reformulation's toolkit is provably not the right shape
to close. The reformulation closes the covering obstruction exactly
where classification suffices, and the conjecture-level consequence is
the one stated in the introduction: within this framework the residual
is a structural gap, not a technical one, and closing it is work for a
different instrument, not another session of this one.

The disposition of that verdict is now on record, and it fixes the
program's final state. The natural external instrument for the missing
margin---the Ehrhart theory of the Lonely-Runner zonotope, the third
bridge of \S\ref{sec:open}---was engaged as a full strand and executed
to certificates. Its arithmetic is unimpeachable: the $h^*$-polynomials
of both canonical families are real-rooted, log-concave, and Newton
through $n=10$, in exact rational arithmetic throughout. Its geometric
reading is false: the deepest-hole form fails from $n=5$ with exact
witnesses (the harmonic lower bound improved to $\rho\ge 4183/12288$),
and the modular law it suggested---deepest hole at the center exactly
when $n\mid m$---holds exactly at $n=3$ ($m\le 24$), is certified at
$n=4$ (through $m=16$, value $9/34$), breaks at $n=5$ (the harmonic
$m=5$ and $m=15$ refuted exactly, $\rho\ge 97/288$ and $6/19$; the
surviving multiples $m=10,20,25,30$ certified at $7/22$, $13/42$,
$4/13$, and $19/62$ respectively---$m=20,25,30$ in the first
certification session and $m=10$ closed by the companion paper's
third-revision attempt, which paired a heavily scaled random-local
search with an exact Lipschitz box certificate and emptied the tree),
and is
extinct at $n=6$ (every tested instance refuted exactly, through the
$4n$ dilate, the narrowest margin $71/208$ against $17/50$). The
instrument that most faithfully preserves the instance's arithmetic is
again the one that cannot produce the margin: the type mismatch is a
property of the lossless encodings themselves, not of one toolkit.
The record is accordingly published in two parts---this monograph (the
exact identities, the classifications, the conditional theorem with
its ledger) and a companion paper \cite{companion-paper} that works
out the negative half: the four falsified strengthenings, the witness
archive, the certification ladder above with its negative control, and
the $h^*$ addendum. The scope statement of the abstract applies to
both parts unchanged.
